% Ссылка в тексте: (приложение~Ж)
\chapter{ТЕСТОВЫЙ КОРПУС ДОКУМЕНТОВ}

В~приложении описан измеренный корпус, использованный в артефакте
\texttt{baseline\_comparison.json} для baseline-прогона B2. Корпус
содержит 19 документов в трёх категориях: продуктово-редакционные
документы, корпоративные регламенты и учебные программы MIT~OCW.

\section*{Описание измеренного корпуса B2}

\begin{table}[h!]
  \centering
  \small
  \caption{Состав измеренного корпуса baseline B2}
  \label{tab:dataset}
  \begin{tabular}{|p{4.0cm}|p{1.5cm}|p{1.5cm}|p{5.4cm}|}
    \hline
    \textbf{Группа} & \textbf{Кол-во} & \textbf{Язык} & \textbf{Назначение в проверке} \\
    \hline
    Продуктово-редакционные документы & 7 & RU &
    Проверка генерации по брендбукам, редакционным правилам и продуктовым описаниям. \\
    \hline
    Корпоративные регламенты & 6 & RU &
    Проверка извлечения требований, ограничений и нормативных формулировок. \\
    \hline
    Учебные программы MIT OCW & 6 & EN &
    Проверка переноса структуры дисциплин и учебных результатов в курс. \\
    \hline
    Итого & 19 & RU/EN &
    Единый baseline-корпус для одношаговой конфигурации B2. \\
    \hline
  \end{tabular}
\end{table}

Machine-readable manifest прогона хранится в
\texttt{baseline\_comparison.json}. В~него входят идентификаторы:
\begin{itemize}
  \item product-summary, VK brand guidelines, гайды по маскоту и мерчу,
  редакционные политики Сбер/HR;
  \item information security policy, personal data policy, incident
  management, onboarding, anticorruption, regulations summary;
  \item MIT OCW 14.01, 18.01, 6.0001, 6.042J, 8.01 и summary-файл.
\end{itemize}

Корпус следует трактовать как технический baseline для проверки
одношаговой конфигурации B2, а не как окончательную экспертно
размеченную выборку. Его задача --- воспроизводимо показать разрыв
между скоростью одношаговой генерации и требованиями к source-grounded
методическому качеству.

\section*{Оценочные артефакты текущей итерации}

На~текущем этапе полностью подготовлены:
\begin{itemize}
  \item измеренный корпус baseline B2 из 19 документов;
  \item набор из~30 вопросно-ответных пар для~RAG-оценки;
  \item набор из~100 утверждений для~оценки факт-чекинга.
\end{itemize}

Paired B2/B4-прогон выполнен на том же наборе документов:
для каждого \texttt{doc\_i} сохранены \texttt{B2\_result\_i},
\texttt{B4\_result\_i}, отдельный \texttt{kb\_slug} B4 и строковая
дельта Bloom/source-grounding. Сводка и сырые REST-артефакты приведены
в~приложении~И и файле \texttt{paired\_b2\_b4\_experiment.json}.
Дополнительно выполнен детерминированный evaluation lab:
\texttt{evaluation\_lab/ai\_methodist\_evaluation\_lab\_latest.json}
фиксирует RAG-style метрики на~30 QA-парах и fact-check
precision / recall / F1 на~100 утверждениях.
